\chapter{Complexity — Meyer Verifier}
%%% generated by the blueprint pipeline from TCSlib.Complexity.CircuitComplexity.MeyerVerifier
\section{Overview}

Declarations of the Lean module \texttt{TCSlib.Complexity.CircuitComplexity.MeyerVerifier}.

\section{Declarations}

\begin{definition}[Atomic formula types]
\lean{Complexity.Meyer.Atom}
\label{Complexity.Meyer.Atom}
\leanok
\uses{Complexity.Meyer.AW, Complexity.Meyer.TW, Complexity.Meyer.TabSel}
For a type $S$ and a natural number $k$, define the type of atoms to consist of an atom
carrying a table selector from $S$ with parameter $k$, a work-tape index, and an
auxiliary work-tape index, together with the distinct atoms for length, the two all-zero
tests, the auxiliary-zero test, the index-zero test, the comparisons of the index with
$N$, the test that the index equals $N+1$, and the input-bit test.
\end{definition}

\begin{theorem}[Membership in a language complement]
\lean{Complexity.Meyer.lang_mem_compl}
\label{Complexity.Meyer.lang_mem_compl}
\leanok
\difficulty{1}
Let $L$ be a language of finite binary strings and let $z$ be a finite binary string.
Then $z$ belongs to the complement of $L$ if and only if $z$ does not belong to $L$.
\end{theorem}

\begin{theorem}[Membership in a predicate-defined language]
\lean{Complexity.Meyer.lang_mem_setOf}
\label{Complexity.Meyer.lang_mem_setOf}
\leanok
\difficulty{1}
Let $p$ be a predicate on finite binary strings and let $z$ be a finite binary string.
Then $z$ belongs to the language consisting of strings satisfying $p$ if and only if
$p(z)$ holds.
\end{theorem}

\begin{definition}[Circuit satisfiability query]
\lean{Complexity.Meyer.ansD}
\label{Complexity.Meyer.ansD}
\leanok
\uses{BoolCircuit.CVAL, Complexity.Meyer.Sel, Complexity.Meyer.encodeSel, Complexity.Meyer.query, Turing.pairEncode}
Given a deterministic finite-state machine $M$ over the binary alphabet, a binary string
$D$, a selector $	ext{sel}$, and binary strings $T$, $A$, and $x$, define the answer to
be true exactly when the pair consisting of $D$ and the query formed from $	ext{sel}$,
$T$, $A$, and $x$ belongs to the language of satisfiable fan-in-two circuits.
\end{definition}

\begin{definition}[Evaluation of an atom]
\lean{Complexity.Meyer.evalAtom}
\label{Complexity.Meyer.evalAtom}
\leanok
\uses{BoolCircuit.CVAL, Complexity.Meyer.Atom, Complexity.Meyer.ansD, Complexity.Meyer.awd, Complexity.Meyer.bjw, Complexity.Meyer.idxDesc, Complexity.Meyer.twd, Complexity.Meyer.vD, Complexity.Meyer.va, Complexity.Meyer.vj, Complexity.Meyer.vt, Complexity.Meyer.vx, Complexity.Meyer.wd, Turing.incFixed, Turing.pairEncode}
Given a deterministic finite-state machine $M$ over the binary alphabet, natural-number
parameters $C_w,c_w$, an atom $a$ associated with $M$, and a binary string $z$, define a
Boolean value by the following cases. For an answer atom, test circuit satisfiability of
the pair formed from the decoded description in $z$ and the query using the atom's
selector and the transformed strings determined by its tags. For the other atoms, test,
respectively: the specified lengths of the decoded strings, including the bound on the
index length; whether incrementing the decoded time or auxiliary word succeeds; whether
the auxiliary word is all zero; whether the decoded index length is zero, less than or
at most the input length, or exactly one more than the input length; or whether the pair
formed from the index description and input belongs to the language of satisfiable
fan-in-two circuits.
\end{definition}

\begin{definition}[Filtering true bits]
\lean{Complexity.Meyer.trues}
\label{Complexity.Meyer.trues}
\leanok
\uses{Complexity.transduce}
For a binary string $u$, define its output to be the string obtained by retaining
exactly the occurrences of $1$ in $u$, in their original order.
\end{definition}

\begin{theorem}[Characterization of empty transducer output]
\lean{Complexity.Meyer.trues_eq_nil}
\label{Complexity.Meyer.trues_eq_nil}
\leanok
\uses{Complexity.Meyer.trues, Complexity.transduce}
\difficulty{4}
For any finite Boolean string $u$, the output produced by the transducer that emits a
true bit exactly when it reads a true bit is empty if and only if every bit of $u$ is
false.
\end{theorem}

\begin{theorem}[Polynomial-time decidability of output-length equality]
\lean{Complexity.Meyer.lenEq_test}
\label{Complexity.Meyer.lenEq_test}
\leanok
\uses{Complexity.P, Complexity.PolyTimeComputable, Complexity.lenEq_preimage_mem_P}
\difficulty{1}
Let $f,g$ be polynomial-time computable functions from finite binary strings to finite
binary strings. The language of strings $z$ for which $|f(z)|=|g(z)|$ belongs to
$\mathsf{P}$.
\end{theorem}

\begin{theorem}[Polynomial-time decidability of output-length comparison]
\lean{Complexity.Meyer.lenLe_test}
\label{Complexity.Meyer.lenLe_test}
\leanok
\uses{Complexity.P, Complexity.PolyTimeComputable, Complexity.lenLe_preimage_mem_P}
\difficulty{1}
Let $f,g$ be polynomial-time computable functions from finite binary strings to finite
binary strings. The language of strings $z$ for which $|g(z)|\le |f(z)|$ belongs to
$\mathsf{P}$.
\end{theorem}

\begin{theorem}[Atomic evaluation is polynomial-time decidable]
\lean{Complexity.Meyer.evalAtom_mem_P}
\label{Complexity.Meyer.evalAtom_mem_P}
\leanok
\uses{BoolCircuit.CVAL_mem_P, Complexity.Meyer.Atom, Complexity.Meyer.ansD, Complexity.Meyer.encodeSel, Complexity.Meyer.evalAtom, Complexity.Meyer.idxDesc, Complexity.Meyer.lenEq_test, Complexity.Meyer.lenLe_test, Complexity.Meyer.pt_awd, Complexity.Meyer.pt_bjw, Complexity.Meyer.pt_inc, Complexity.Meyer.pt_query, Complexity.Meyer.pt_twd, Complexity.Meyer.pt_vD, Complexity.Meyer.pt_va, Complexity.Meyer.pt_vj, Complexity.Meyer.pt_vt, Complexity.Meyer.pt_vx, Complexity.Meyer.trues, Complexity.Meyer.trues_eq_nil, Complexity.Meyer.va, Complexity.Meyer.vj, Complexity.Meyer.vx, Complexity.Meyer.wd, Complexity.P, Complexity.PolyHierarchy.unaryPT_poly, Complexity.PolyTimeComputable, Complexity.PolyTimeComputable.append, Complexity.PolyTimeComputable.comp, Complexity.PolyTimeComputable.pairEncode, Complexity.compl_mem_P, Complexity.inter_mem_P, Complexity.polyTimeComputable_const, Complexity.polyTimeComputable_prepend, Complexity.polyTimeComputable_transduce, Complexity.polyTimeComputable_unary, Complexity.preimage_mem_P}
\difficulty{5}
Let $C_w,c_w\in\mathbb{N}$, let $M$ be a finite Turing machine over the Boolean
alphabet, and let $a$ be an atom over the states and work tapes of $M$. The language of
finite Boolean words $z$ for which the evaluation of $a$ on $z$ is true belongs to $P$,
the class of languages decidable in polynomial time.
\end{theorem}

\begin{definition}[Finite checking alphabet]
\lean{Complexity.Meyer.Check}
\label{Complexity.Meyer.Check}
\leanok
\uses{Complexity.TimeHierarchy.OutReg}
For a type $S$ and a natural number $k$, the type of checks has constructors for an
initial state and output-register value; an initial tape symbol with two Boolean
annotations; an initial index with two Boolean annotations; a zero-negation check
carrying an optional tape position and a Boolean; a state transition carrying old and
new state/output-register data and the scanned-symbol and tape-update data; a tape-cell
transition check carrying that same transition data together with a position and
Boolean; an index transition check carrying that transition data and a Boolean; a
boundary check carrying the transition data; and a final check.
\end{definition}

\begin{definition}[Symbol-value constraint]
\lean{Complexity.Meyer.symvA}
\label{Complexity.Meyer.symvA}
\leanok
\uses{Complexity.Meyer.AW, Complexity.Meyer.TW}
Given a Boolean-valued function $v$ on atoms, an optional tape position $\tau$, a bit
$s$, a time-word tag $T$, and an auxiliary-word tag $A$, define an optional bit as
follows: if $v$ is true on the answer atom with symbol selector $(\tau,s,1)$ and tags
$T,A$, return no value; otherwise return the value of $v$ on the corresponding answer
atom with selector $(\tau,s,0)$ and the same tags.
\end{definition}

\begin{definition}[State-value constraint]
\lean{Complexity.Meyer.stA}
\label{Complexity.Meyer.stA}
\leanok
\uses{Complexity.Meyer.TW, Complexity.TimeHierarchy.OutReg}
Given a Boolean-valued function $v$ on atoms, a state descriptor consisting of an
optional machine state and an output-register value, and a time-word tag $T$, define a
Boolean value by evaluating $v$ on the answer atom for that state descriptor, time tag,
and the zero auxiliary-word tag.
\end{definition}

\begin{definition}[Transition claims]
\lean{Complexity.Meyer.claims}
\label{Complexity.Meyer.claims}
\leanok
\uses{Complexity.Meyer.Rd, Complexity.Meyer.TW, Complexity.Meyer.stA, Complexity.Meyer.symvA, Complexity.TimeHierarchy.OutReg}
Given a Boolean-valued function $v$ on atoms, a time-word tag $T$, a state descriptor
$\sigma$, and register data $\rho$ consisting of an optional output bit and a partial
tape assignment, define a Boolean value that is true exactly when all of the following
hold: $v$ accepts the answer atom for $\sigma$, $T$, and the zero auxiliary-word tag;
the symbol-value constraint at no tape position with input bit $0$ equals the optional
output bit in $\rho$; and at every tape position, the symbol-value constraint with input
bit $0$ equals the value assigned by $\rho$.
\end{definition}

\begin{definition}[Input-value constraint]
\lean{Complexity.Meyer.xinV}
\label{Complexity.Meyer.xinV}
\leanok
Given a Boolean-valued function $v$ on atoms and a bit $s$, define an optional bit to be
absent if $s=1$ and the zero-index atom is false under $v$; otherwise, if the
index-less-than-input-length atom is true under $v$, return the value of the input-bit
atom, and return no value in all other cases.
\end{definition}

\begin{definition}[Work tape consistency condition]
\lean{Complexity.Meyer.workOK}
\label{Complexity.Meyer.workOK}
\leanok
\uses{Complexity.Meyer.AW, Complexity.Meyer.Rd, Complexity.Meyer.symvA, Complexity.Meyer.wsym, Complexity.TimeHierarchy.OutReg}
Let $M$ be a finite-tape Turing machine with Boolean tape alphabet, let $v$ be a Boolean
valuation of its atoms, and let $\rho$ consist of an optional input symbol and an
optional Boolean read symbol for each work tape. For a state-or-empty value $	au$,
output register $r$, work-tape index $h$, and Boolean $s$, set $c_1=a_n$ and $c_2=a_c$
when $s$ is true, and reverse these choices when $s$ is false. The condition holds if
$	au$ is empty precisely when the two symbols represented by the $t_n$ and $t_c$ answer
atoms at $(h,s,c_1)$ agree; if $	au$ is a state $q$, use the transition from $q$ on
$\rho$ and its move direction: for a positive move, the $t_n$ symbol at $(h,s,c_1)$ must
equal the $t_c$ symbol at $(h,s,c_2)$, except that when $s$ and the zero-answer atom are
true it is compared instead with the updated work-tape symbol (or the read symbol if no
update is specified); for a negative move, the $t_n$ symbol at $(h,s,c_2)$ must equal
the $t_c$ symbol at $(h,s,c_1)$, with that same substitution when $
eg s$ and the zero-answer atom are true; for a stationary move, the $t_n$ and $t_c$
symbols at $(h,s,c_1)$ must agree, with that substitution when $
eg s$ and the zero-answer atom are true.
\end{definition}

\begin{definition}[Clamped input condition]
\lean{Complexity.Meyer.clampedA}
\label{Complexity.Meyer.clampedA}
\leanok
\uses{Complexity.Meyer.Rd, Complexity.Meyer.symvA}
Let $M$ be a finite-tape Turing machine with Boolean tape alphabet, let $v$ be a Boolean
valuation on its atoms, let $\rho$ specify an optional input symbol and the read symbols
on its work tapes, and let $m$ be a sign. The condition holds exactly when the input
symbol in $\rho$ is absent and the symbol encoded by the $t_c$ answer atom of type
$a_1$, with no work tape selected and with Boolean parameter true precisely when $m$ is
negative, is absent.
\end{definition}

\begin{definition}[Input tape consistency condition]
\lean{Complexity.Meyer.inOK}
\label{Complexity.Meyer.inOK}
\leanok
\uses{Complexity.Meyer.AW, Complexity.Meyer.Rd, Complexity.Meyer.clampedA, Complexity.Meyer.symvA, Complexity.TimeHierarchy.OutReg}
Let $M$ be a finite-tape Turing machine with Boolean tape alphabet, let $v$ be a Boolean
valuation on its atoms, let $\sigma$ be an optional machine state paired with an output
register, let $\rho$ specify read data, and fix a Boolean $s$. First require that, for
each of the two answer types $a_j$ and $a_{j1}$ (in the order determined by $s$), the
symbols encoded by the corresponding $t_n$ and $t_c$ answer atoms agree. If $\sigma$ has
no state, or its transition reads the input tape at position zero, this is the whole
condition. If the transition reads a positive or negative position and the input is not
clamped (the input symbol in $\rho$ is absent and the specified $t_c,a_1$ answer symbol
is absent), instead require equality only between the $t_n$ symbol of the first answer
type and the $t_c$ symbol of the second for a positive position, or between the $t_n$
symbol of the second and the $t_c$ symbol of the first for a negative position; if the
input is clamped, retain the initial pair of equalities. The first answer type is
$a_{j1}$ when $s$ is true and $a_j$ otherwise, and the second is the reverse.
\end{definition}

\begin{definition}[Input boundary consistency condition]
\lean{Complexity.Meyer.bdOK}
\label{Complexity.Meyer.bdOK}
\leanok
\uses{Complexity.Meyer.Rd, Complexity.Meyer.clampedA, Complexity.Meyer.symvA, Complexity.TimeHierarchy.OutReg}
Let $M$ be a finite-tape Turing machine with Boolean tape alphabet, let $v$ be a Boolean
valuation on its atoms, let $\sigma$ be an optional machine state paired with an output
register, and let $\rho$ specify read data. The condition is true if $\sigma$ has no
state or the transition from its state reads the input tape at position zero. For a
positive input position, it is true exactly when the input is clamped or the symbol
encoded by the $t_n$ answer atom of type $a_j$, with no work tape selected and Boolean
parameter false, is absent; for a negative position, replace that answer atom's Boolean
parameter by true. The input is clamped when the input symbol in $\rho$ is absent and
the $t_c$ answer symbol of type $a_1$, with no work tape selected and Boolean parameter
true precisely when the input position is negative, is absent.
\end{definition}

\begin{definition}[Local verification predicate]
\lean{Complexity.Meyer.goodF}
\label{Complexity.Meyer.goodF}
\leanok
\uses{Complexity.Meyer.Check, Complexity.Meyer.bdOK, Complexity.Meyer.claims, Complexity.Meyer.inOK, Complexity.Meyer.ruleS, Complexity.Meyer.symBit, Complexity.Meyer.workOK, Complexity.Meyer.xinV}
Let $M$ be a finite-tape Turing machine with Boolean tape alphabet and initial state
$q_0$, and let $v$ be a Boolean valuation on the atoms associated with its states and
tapes. This predicate assigns a Boolean to each verification item as follows: an
initial-state item checks that the state and output register are $(q_0,\mathrm{empty})$
exactly when the corresponding $t_0,a_0$ atom is true; an initial work-tape item checks
that the $t_0,a_c$ atom agrees with whether the represented symbol is blank; and an
initial input item checks the analogous $t_0,a_j$ atom against the input symbol
specified by the valuation's input-length and input-bit atoms. A zero-sign item requires
the two $t_c,a_0$ state answers for opposite sign parameters to agree. A
state-transition item imposes the next-state and output-register update when the time
atoms are all true and the current configuration and read data satisfy their claims; a
work-tape transition item imposes the work-tape consistency condition when those time
atoms and all answer atoms are true and the claims hold; an input-transition item
imposes the input-tape consistency condition when the time and input-index bounds and
claims hold; and an input-boundary item imposes the boundary consistency condition when
the time and input-index boundary atoms and claims hold. These four transition checks
are vacuously true when their indicated conditions do not hold. A final item requires
the $t_1,a_0$ state answer for no state and a one-symbol output register to be true.
\end{definition}

\begin{definition}[Bad valuation condition]
\lean{Complexity.Meyer.badF}
\label{Complexity.Meyer.badF}
\leanok
\uses{Complexity.Meyer.Check, Complexity.Meyer.goodF}
Let $M$ be a finite-tape Turing machine with Boolean tape alphabet, and let $v$ be a
Boolean valuation on its atoms. The condition holds exactly when the valuation's
length-all atom is true and there exists a verification item on which the local
verification predicate is false.
\end{definition}

\begin{definition}[Language of bad encodings]
\lean{Complexity.Meyer.badLang}
\label{Complexity.Meyer.badLang}
\leanok
\uses{Complexity.Meyer.badF, Complexity.Meyer.evalAtom}
Let $C_w,c_w$ be natural numbers and let $M$ be a finite-tape Turing machine with
Boolean tape alphabet. The language consists of precisely those Boolean words $z$ for
which the valuation of machine atoms obtained by evaluating them with parameters
$C_w,c_w,M$ on $z$ satisfies the bad-valuation condition.
\end{definition}

\begin{theorem}[Polynomial-time decidability of the bad language]
\lean{Complexity.Meyer.badLang_mem_P}
\label{Complexity.Meyer.badLang_mem_P}
\leanok
\uses{Complexity.Meyer.badF, Complexity.Meyer.badLang, Complexity.Meyer.evalAtom, Complexity.Meyer.evalAtom_mem_P, Complexity.P, Complexity.mem_P_of_atoms}
\difficulty{1}
For any natural numbers $C_w,c_w$ and any finite-state Turing machine $M$ with Boolean
alphabet, the language defined by the bad-valuation predicate for these parameters
belongs to polynomial time.
\end{theorem}

\begin{theorem}[Congruence of atomic evaluations]
\lean{Complexity.Meyer.evalAtom_congr}
\label{Complexity.Meyer.evalAtom_congr}
\leanok
\uses{Complexity.Meyer.Atom, Complexity.Meyer.awd, Complexity.Meyer.bjw, Complexity.Meyer.evalAtom, Complexity.Meyer.twd, Complexity.Meyer.vD, Complexity.Meyer.va, Complexity.Meyer.vj, Complexity.Meyer.vp, Complexity.Meyer.vt, Complexity.Meyer.vx}
\difficulty{4}
Let $C_w,c_w\in\mathbb{N}$, let $M$ be a finite Turing machine over the Boolean
alphabet, and let $z,z'$ be finite Boolean words. Suppose their six decoded
components—representing the input, description, time word, auxiliary word, index word,
and final word—agree componentwise. Then every atom over the states and work tapes of
$M$ has the same Boolean evaluation on $z$ and $z'$.
\end{theorem}

\begin{definition}[Inner language]
\lean{Complexity.Meyer.innerLang}
\label{Complexity.Meyer.innerLang}
\leanok
\uses{Complexity.Meyer.badLang, Turing.pairEncode}
Let $C_w,c_w$ be natural numbers and let $M$ be a finite-tape Turing machine with
Boolean tape alphabet. The inner language consists of exactly those Boolean words $w$
such that, for every Boolean word $v$, the paired encoding of $w$ and $v$ does not
belong to the language of bad encodings with parameters $C_w,c_w$ and machine $M$.
\end{definition}

\begin{theorem}[CoNP membership of the inner language]
\lean{Complexity.Meyer.innerLang_mem_coNP}
\label{Complexity.Meyer.innerLang_mem_coNP}
\leanok
\uses{Complexity.Meyer.badF, Complexity.Meyer.badLang, Complexity.Meyer.badLang_mem_P, Complexity.Meyer.bjw, Complexity.Meyer.evalAtom, Complexity.Meyer.evalAtom_congr, Complexity.Meyer.innerLang, Complexity.Meyer.vD, Complexity.Meyer.va, Complexity.Meyer.vj, Complexity.Meyer.vp, Complexity.Meyer.vt, Complexity.Meyer.vx, Complexity.Meyer.wd, Complexity.NP, Complexity.coNP, Complexity.length_pairFstD_le, Complexity.mem_NP_iff_exists_length_le, Complexity.pairFstD_pairEncode, Turing.length_pairEncode, Turing.pairEncode}
\difficulty{5}
Let $C_w,c_w\in\mathbb{N}$ and let $M$ be a finite Turing machine over the Boolean
alphabet. The language of Boolean strings $w$ for which every Boolean string $v$ has the
property that the pair encoding of $w$ and $v$ does not belong to the language of
malformed or locally inconsistent encodings defined by the associated verifier is in
coNP.
\end{theorem}
